% 第1章 动态宏观经济学的基本方法与概念

% 导言

\chapter{最优增长与动态规划基础} \label{ch-ngm-dp}

动态规划是量化宏观经济学的基础。本章借助最优增长模型这一现代宏观经济学的基石模型（Workhouse Model），
通过将原始的序列问题转化为递归形式问题，进而引入确定性动态规划的基本理论方法，建立了贝尔曼方程，对其解的存在性和相应性质进行分析，
并介绍了这一模型的初步数值解法，为后续学习奠定基础。


% ===== 1.1 确定性动态规划 =====
\section{最优增长模型} \label{sec-ngm}

现代增长理论始于Ramsey（1928）的经典论文，此后Solow（1956）和Swan（1956）的工作标志这新古典增长理论的诞生。

首先回顾经典的Solow模型。总量生产函数：
\begin{equation*}
    Y_t = F\left(K_t, L_t\right)
\end{equation*}

Solow模型的主要问题在于，模型中储蓄率是外生给定的。Cass（1965）和Koopmans（1965）在Solow工作的基础上，对新古典理论进行了拓展，
在模型中引入家庭偏好，同时将资本积累内生化，建立了确定性情形下的最优增长理论。因此最优增长模型通常被称为RCK模型（Ramsey-Cass-Koopmans Model）。

典型的确定性最优增长模型形式如下：
假设时间是离散且无穷的，由$t=0,1,\dots, \infty$表示，给定第0期资本存量：
$$k_0>0$$
中央计划者（Social Planner）在资源约束:
$$c_t+k_{t+1}=f\left(k_t\right)+(1-\delta)k_t,\quad \forall t \geq 0$$
和资源非负的约束下：
$$c_t \geq 0, \quad k_{t+1} \geq 0, \quad \forall t \geq 0$$
选择无穷期序列：
$$\left\{c_t, k_{t+1}\right\}_{t=0}^{\infty}$$
最大化无穷期效用函数：
$$\sum_{t=0}^{\infty} \beta^t u\left(c_t\right)$$

这一形式通常被称为\textbf{序列问题（Sequence Probem）}。同时，模型还有如下的规范化假设，我们将在下一部分具体理解这些假设的用途。
\begin{assumption} \label{as-f}
    生产函数$f\left(k\right)$连续，二次可微
        且$f\left(0\right)=0,\quad f^{\prime}\left(k\right)>0,\quad f^{\prime\prime}\left(k\right)<0$
\end{assumption}
\begin{assumption} \label{as-f-inada}
    生产函数Inada条件：
        $\lim\limits_{k \rightarrow 0} f^{\prime}(k)=\infty, \quad \lim\limits_{k \rightarrow \infty} f^{\prime}(k)=0$
\end{assumption}
\begin{assumption} \label{as-u}
    效用函数$u$连续，二次可微
        且$u^{\prime}\left(c\right)>0,\quad u^{\prime\prime}\left(c\right)<0$
\end{assumption}
\begin{assumption} \label{as-u-inada}
    效用函数Inada条件：
    $\lim\limits_{c \rightarrow 0} u^{\prime}(c)=\infty, \quad \lim\limits_{c \rightarrow \infty} U^{\prime}(c)=0$
\end{assumption}


% ===== 1.2 序列问题：基本Lagrangian方法与横截性条件 =====
\section{序列问题：基本Lagrangian解法与横截性条件} \label{sec-sp}

利用Lagrangian函数是处理优化问题的基本方法。但当问题拓展到上面无穷期形式时，需要一些额外条件来确保问题的可解性。我们将从有限期情形开始，逐步推广至无限期的情形。

% 1.2.1 有限期情形
\subsection{有限期情形} \label{subsec-finite-h}

首先考虑上面问题的有限期形式，
即时间$t \in$ $\{0,1, \ldots T\}$。对于最大化问题，记第$t$期资源约束相关的乘子为$\beta^t \lambda_t$，进而可构造如下拉格朗日方程：
$$
\mathcal{L}=\sum_{t=0}^T \beta^t u\left(c_t\right)-\sum_{t=0}^T \beta^t \lambda_t\left[c_t+k_{t+1}-(1-\delta) k_t-f\left(k_t\right)\right]
$$
也即：
$$
\mathcal{L}=\sum_{t=0}^T \beta^t\left\{u\left(c_t\right)-\lambda_t\left[c_t+k_{t+1}-(1-\delta) k_t-f\left(k_t\right)\right]\right\}
$$
不难看出，$\lambda_t$是计划者第$t$期资源约束的影子价格，或者说是边际价值。在生产函数和效用函数满足规范化假设时，只要$k_t>0$，这一问题有内点解。
分别对$c_t$和$k_{t+1}$求一阶条件，可得：
$$u^{\prime}\left(c_t\right)-\lambda_t=0$$
$$-\lambda_t+\beta\lambda_{t+1}\left[1-\delta+f^{\prime}\left(k_{t+1}\right)\right]=0$$
与乘子$\lambda_t$相关的一阶条件即为资源约束：在第$t$期，如果经济额外获得$\epsilon$足够小单位的商品，那么从第$t$期看，
福利增加了$\lambda_t\epsilon$；从第0期看，则福利增加了$\beta^t\lambda_t$单位。
合并两个一阶条件可得：
$$
\frac{u^{\prime}\left(c_t\right)}{\beta u^{\prime}\left(c_{t+1}\right)}=f^{\prime}\left(k_{t+1}\right)+1-\delta
$$
这一条件表明：计划者的最优选择是使得今明两天消费的MRS与MRT相等

对于资本$k$，一个自然的问题是：
终期$T$时如何选择，即$k_{T+1}$的最优值?由$k_{T+1}$的KT条件:
$$\frac{\partial \mathcal{L}}{\partial k_{T+1}} \geq 0, \quad k_{T+1} \geq 0$$
且满足互补松弛条件。这等价于：
$$
\beta^T \lambda_T \geq 0, \quad k_{T+1} \geq 0, \quad \beta^T \lambda_T k_{T+1}=0
$$
这意味着在终期第$T$期选择的$k_{T+1}=0$或$k_{T+1}$的影子价值为0，因为不存在$T+1$期，从而在第$T$期选择$k_{T+1}$没有意义。
又由于$\lambda_T=u^{\prime}\left(c_T\right)$，而模型\textbf{假设}$u^{\prime} > 0$，因此，终期T存在角点解：
$$k_{T+1}=0$$
或者说，如果终期选择$k_{T+1}>0$，则计划者会余下未使用资源，在规范化假设下这不会是最优策略（因为$u^{\prime}>0$）。

% 1.2.2 无限期情形
\subsection{无限期情形} \label{subsec-infinit-h}

现在，\textbf{由有限期转向无限期问题}。回忆在有限期问题中，终期时$k_{T+1}$的最优条件满足：
$$\beta^T \lambda_T k_{T+1}=0$$
从直觉出发，令$T \rightarrow \infty$，那么可以得到无限期情形下的\textbf{横截性条件}（Transversality Condition）：
$$
\lim _{T \rightarrow \infty} \beta^T \lambda_T k_{T+1}=0
$$

无限期最优增长模型中，计划者的最优配置$\left\{c_t, k_{t+1}\right\}_{t=0}^{\infty}$是符合
如下欧拉方程（optimality）、资源约束（feasibility）和边界条件的一条\textbf{唯一路径}（unique path）：
\begin{equation}    \label{eq-euler-sp}
    \frac{u^{\prime}\left(c_t\right)}{\beta u^{\prime}\left(c_{t+1}\right)}=f^{\prime}\left(k_{t+1}\right)+1-\delta
\end{equation}
\begin{equation}    \label{eq-res-const-sp}
    k_{t+1}=f\left(k_t\right)+(1-\delta) k_t-c_t
\end{equation}
以及两个边界条件
$$
k_0>0, \quad \lim _{t \rightarrow \infty} \beta^t u^{\prime}\left(c_t\right) k_{t+1}=0
$$
式\eqref{eq-euler-sp}被称为\textbf{欧拉方程（Euler Equation）}。
事实上，符合前两条条件的$\left\{c_t, k_{t+1}\right\}_{t=0}^{\infty}$路径可能有多条，但在两个边界条件约束下，这条路径是唯一确定的。后续将会证明。


% 1.3 ===== 序列问题到递归问题 =====
\section{动态规划：由序列问题转向递归形式}  \label{sec-sp-fe}

前一节特意强调是在求解一个\textbf{序列问题}，优化后得到的解是一些无穷序列数值。但这一形式的求解存在一定不足：
处理无限期界过于复杂，需要求解无穷的差分方程序列，我们希望降维。不难发现，当计划者在第0期选择了$c_0,k_1$后，其余选择可以等到下一期做出。更一般的来说，
无穷序列问题可以分解为一期一期的小问题：对于$t+1$期，可以在第$t$期进行选择而不用在第0期进行，即定义从当期到下期的映射，
这一问题形式被称为\textbf{递归形式}（Recursive Formulation）。
\textbf{序列问题求解后，我们得到了无穷个数值（Number）；对于递归问题，我们求解所得的是一个方程（Function）,其本质是一个泛函问题}。
利用泛函方程（Function Equation）对动态优化问题进行研究就是\textbf{动态规划}（Dynamic Programmming）。
动态规划最早由Richard E. Bellman提出，是现代宏观经济学的核心。


% 1.3.1 值函数、状态变量与贝尔曼方程
\subsection{概念：值函数、状态变量与贝尔曼方程} \label{subsec-dp-concept}

对于最优增长问题，唯一可变变量是$k_0$：给定不同的初始资本将会得到不同的最优化序列和最大化效用值。因此，可以定义最大化目标效用值的函数：
$$
v\left(k_0\right) \equiv \max _{\left\{k_{t+1}\right\}_{t=0}^{\infty}} \sum_{t=0}^{\infty} \beta^t u\left(f\left(k_t\right)-k_{t+1}\right), \quad 0<\beta<1
$$
这一函数称为值函数（Value Function）。记$\gamma\left(k_t\right)\equiv f\left(k_t\right)+\left(1-\delta\right)k_t$，为了看清递归形式的特征，对这一求和进行如下变换：
$$
\begin{aligned}
v\left(k_0\right) & \equiv \max _{\left\{k_{t+1}\right\}_{t=0}^{\infty}} \sum_{t=0}^{\infty} \beta^t u\left(\gamma\left(k_t\right)-k_{t+1}\right) \\
& =\max _{\left\{k_{t+1}\right\}_{t=0}^{\infty}}\left[u\left(\gamma\left(k_0\right)-k_1\right)+
    \beta \sum_{t=1}^{\infty} \beta^{t-1} u\left(\gamma\left(k_t\right)-k_{t+1}\right)\right] \\
& =\max _{k_1}\left[u\left(\gamma\left(k_0\right)-k_1\right)+
    \beta \max _{\left\{k_{t+1}\right\}_{t=1}^{\infty}} \sum_{t=1}^{\infty} \beta^{t-1} u\left(\gamma\left(k_t\right)-k_{t+1}\right)\right] \\
& =\max _{k_1}\left[u\left(\gamma\left(k_0\right)-k_1\right)+\beta v\left(k_1\right)\right]
\end{aligned}
$$
由此，第0期的一个无穷序列问题转化成了一个具有递归特征的问题：期初值函数（目标）$v\left(k_0\right)$转化成了当期效用+下期值函数折现的\textbf{递归形式}。
由于时间是无穷的，从第0期开始和从第1期开始并没有本质区别，因此递归问题可以进一步写为：
$$
v\left(k_t\right)=\max _{k_{t+1}}\left[u\left(\gamma\left(k_t\right)-k_{t+1}\right)+\beta v\left(k_{t+1}\right)\right]
$$

一个遗留的问题是：值函数$v$是否会依赖时间$t$？答案是，只要原问题是平稳的（stationary），值函数$v$就不会依赖于时间。直觉上看，
虽然状态变量的值（值函数的自变量）可能不同，但只要原问题在第$t$和第$t+1$期之间是相似的，值函数形式本身就不会依赖于时间。但并非所有问题都是平稳的，
例如，在有限期的最优增长模型中，计划者选择资本时会关注剩余的时间。随着终期的接近，决策问题也会发生变化。此时，问题虽然仍可表征为递归形式，但值函数会依赖于时间，即$v_t$。
对于无限期最优增长模型，由于问题具有平稳性，因此递归形式中时间记号完全不重要，值函数被一般化的记为：
\begin{equation}    \label{eq-bellman}
    v\left(k\right)=\max _{k^{\prime}}\left[u\left(\gamma\left(k\right)-k^{\prime}\right)+\beta v\left(k^{\prime}\right)\right]
\end{equation}
这一方程被称为\textbf{贝尔曼方程（Bellman Equation）}。更一般的，这是一个\textbf{泛函方程（Functional Equation）}。因为到目前为止我们仍不知道值函数$v$的形式，
因此下面的目标就是求解这一泛函问题（求解一个函数）。更具体来说，我们的目标是，给定函数$u$和$f$的外生形式下，求解内生值函数$v$。

不难看出，这一问题中对计划者当期决策唯一有用的信息是当期资本存量$k_t$，这样的变量被称为\textbf{状态变量}（State Variable）。
由于这一变量在$t-1$期决策时就已被确定，也被称为\textbf{前定变量}（Predetermined Variable）。


% 1.3.2 FOC Euler
\subsection{一阶条件、策略函数、包络定理与欧拉泛函} \label{subsec-func-euler}


% 1.3.2.1 由一阶条件到欧拉泛函
\subsubsection{由一阶条件到欧拉泛函}

对于贝尔曼方程\ref{eq-bellman}：
$$
v\left(k\right)=\max _{k^{\prime}}\left[u\left(\gamma\left(k\right)-k^{\prime}\right)+\beta v\left(k^{\prime}\right)\right]
$$
观察方程右边，\textbf{假设值函数$v\left(\cdot\right)$已知}，那么就简化为一个两期优化问题：
$$
\max _{k^{\prime}}\left[u\left(\gamma\left(k\right)-k^{\prime}\right)+\beta v\left(k^{\prime}\right)\right]
$$
对$k^{\prime}$求导可得\textbf{一阶条件}：
\begin{equation}    \label{bellman-foc}
    u^{\prime}\left(\gamma\left(k\right)-k^{\prime}\right) = \beta v^{\prime}\left(k^{\prime}\right)
\end{equation}
我们希望从中求解出某一\textbf{策略函数}（Policy Function），或称为\textbf{决策规则}（Decision Rule）：
$k^{\prime}=g\left(k\right)$：
$$
g\left(k\right) \equiv \underset{k^{\prime}}{\operatorname{argmax}}[u(\gamma\left(k\right)-k^{\prime})+\beta v\left(k^{\prime}\right)]
$$
但\textbf{问题在于}：$v$未知，进而$v^{\prime}$也未知。

在问题取到最优$k^{\prime}=g\left(k\right)$时，Bellman方程的最大化算子消去，且对任意$k$成立
$$
v\left(k\right)=u\left(\gamma\left(k\right)-g\left(k\right)\right)+\beta v\left(g\left(k\right)\right),\quad \forall k
$$
为了求解$v^{\prime}$，方程两边同时对$k$做微分可得：
\begin{equation}
    \begin{aligned}
    v^{\prime}\left(k\right) 
    &= u^{\prime}\left(\gamma\left(k\right)-g\left(k\right)\right)\left(\gamma^{\prime}\left(k\right)-g^{\prime}\left(k\right)\right)
        + \beta v^{\prime}\left(g\left(k\right)\right)g^{\prime}\left(k\right) \\
    &= u^{\prime}\left(\gamma\left(k\right)-g\left(k\right)\right)\gamma^{\prime}\left(k\right) 
        -u^{\prime}\left(\gamma\left(k\right)-g\left(k\right)\right)g^{\prime}\left(k\right) 
        + \beta v^{\prime}\left(g\left(k\right)\right)g^{\prime}\left(k\right) \\
    & = u^{\prime}\left(\gamma\left(k\right)-g\left(k\right)\right)\gamma^{\prime}\left(k\right)
        +g^{\prime}\left(k\right)\left[-u^{\prime}\left(\gamma\left(k\right)-g\left(k\right)\right)+\beta v^{\prime}\left(g\left(k\right)\right)\right] \\
    &= \underbrace{u^{\prime}\left(c\right)\left[f^{\prime}\left(k\right)+\left(1-\delta\right)\right]}_{\text{k变化的直接效应}}
            + \underbrace{\left[-u^{\prime}\left(c\right)+\beta v^{\prime}\left(g\left(k\right)\right)\right]g^{\prime}\left(k\right)}
                _{\text{选择最优$k^{\prime}=g\left(k\right)$的间接效应}}
    \end{aligned}
\end{equation}
由于策略函数满足一阶条件，即：
$$
u^{\prime}\left(\gamma\left(k\right)-g\left(k\right)\right)=\beta v^{\prime}\left(g\left(k\right)\right)
$$
将其带入全微分的恒等式,发现右侧第二项——也即由选择最优$k^{\prime}=g(k)$带来的间接影响为0，进而就得到值函数的导数形式：
\begin{align}
    v^{\prime}\left(k\right) &= u^{\prime}\left(\gamma\left(k\right)-g\left(k\right)\right)\gamma^{\prime}\left(k\right) \nonumber \\
        &= u^{\prime}\left(c\right)\left[f^{\prime}\left(k\right)+1-\delta\right] \label{bs}
\end{align}
不难看出，推导过程实际就是\textbf{包络定理}的应用：状态变量$k$的变化会通过两条途径影响值函数：
\begin{itemize}
    \item 直接路径：$k$变化-当期产出$f\left(k\right)$变化-消费$c$变化
    \item 间接路径：$k$变化-最优选择$k^{\prime}$变化-影响未来效用$v\left(k^{\prime}\right)$
\end{itemize}
包络定理表明，由于一阶条件的最优性，间接效应被抵消，因此只剩下直接效应。在动态规划问题中，上述条件也被称为\textbf{Benveniste-Schenkman条件}

基于所得到的值函数导数$v^{\prime}\left(k\right)$公式\eqref{bs}，容易得到：
$$
v^{\prime}\left(k^{\prime}\right)
    =u^{\prime}\left(\gamma\left(k^{\prime}\right)-g\left(k^{\prime}\right)\right)\gamma^{\prime}\left(k^{\prime}\right)
$$
进而可以将一阶条件\eqref{bellman-foc}进一步写为：
\begin{align}
    u^{\prime}\left(\gamma\left(k\right)-g\left(k\right)\right) 
        &= \beta v^{\prime}\left(g\left(k\right)\right) \nonumber    \\
        &= \beta u^{\prime}\left(\gamma\left(g\left(k\right)\right)-g\left(g\left(k\right)\right)\right)
            \gamma^{\prime}\left(g\left(k\right)\right) \label{func-euler}
\end{align}
即：
\begin{equation*}
    u^{\prime}\left(c\right)=\beta u^{\prime}\left(c^{\prime}\right)\left[f^{\prime}\left(k^{\prime}\right)+\left(1-\delta\right)\right]
\end{equation*}
式\eqref{func-euler}就是\textbf{欧拉方程（Euler Equation）}。回顾此前序列问题中得到的欧拉方程：
$$
\frac{u^{\prime}\left(c_t\right)}{\beta u^{\prime}\left(c_{t+1}\right)}=f^{\prime}\left(k_{t+1}\right)+1-\delta
$$
实际也符合这一形式，但其中的重要区别在于：序列问题的最优方程中，要求解的未知量是资本的无穷序列值；
但递归问题的欧拉方程\eqref{func-euler}本质上是一个泛函，对所有$k$成立，\textbf{未知的是策略函数}$k^{\prime}=g\left(k\right)$。
在\textbf{递归形式}下，欧拉方程称为欧拉泛函（Functional Euler Equation）。

% 1.3.2.2 Bellman vs Euler

\subsubsection{贝尔曼方程 vs 欧拉方程}
在对贝尔曼方程进行研究时，我们提到希望从中求解出\textbf{策略函数}（Policy Function）
$$
k^{\prime}_{t+1}=g\left(k_t\right), \quad k_0>0
$$
这是关于$k_t$的一维动力系统（Dynamic System）；
而如果直接从欧拉方程的角度看：
$$u^{\prime}\left(\gamma\left(k_t\right)-k_{t+1}\right) 
        = \beta u^{\prime}\left(\gamma\left(k_{t+1}\right)-k_{t+2}\right)
            \gamma^{\prime}\left(k_{t+1}\right)
$$
这是一个关于$k_t$的二维动力系统（Dynamic System）。为何存在差异？
实际上，这对应了两种不同解法。回顾贝尔曼方程和欧拉泛函：
$$
v\left(k\right)=\max _{k^{\prime}}\left[u\left(\gamma\left(k\right)-k^{\prime}\right)+\beta v\left(k^{\prime}\right)\right]
$$
$$u^{\prime}\left(\gamma\left(k\right)-g\left(k\right)\right) 
        = \beta u^{\prime}\left(\gamma\left(g\left(k\right)\right)-g\left(g\left(k\right)\right)\right)
            \gamma^{\prime}\left(g\left(k\right)\right)
$$
二者都是在求解泛函问题，区别在于：
\begin{itemize}
    \item 贝尔曼方程的核心目标是值函数$v$，在得到值函数的基础上确定策略函数$g\left(k\right)$
    \item 欧拉方程的目标是直接从泛函中解出略函数$g\left(k\right)$
\end{itemize}
在此后将会看到，前者更稳健，后者速度更快。

至此，进行一个小结。我们已经将求解无限期的序列问题（Sequential Problem）转换为了
递归形式（Recursive Formulation）的泛函问题：首先建立了贝尔曼方程和一阶条件，定义了策略函数；
接着利用包络定理（BS-Theorm）得到欧拉泛函。
\begin{itemize}
    \item Bellman Equation：
    $v\left(k\right) = \max _{k^{\prime}}\left[u\left(\gamma\left(k\right)-k^{\prime}\right)+\beta v\left(k^{\prime}\right)\right]$
    \item FOC：$u^{\prime}\left(\gamma\left(k\right)-g\left(k\right)\right) = \beta v^{\prime}\left(g\left(k\right)\right)$
    \item 包络定理：$v^{\prime}\left(k\right) = u^{\prime}\left(\gamma\left(k\right)-g\left(k\right)\right)\gamma^{\prime}\left(k\right)$
    \item 欧拉泛函：$u^{\prime}\left(\gamma\left(k\right)-g\left(k\right)\right) 
        = \beta u^{\prime}\left(\gamma\left(g\left(k\right)\right)-g\left(g\left(k\right)\right)\right)
            \gamma^{\prime}\left(g\left(k\right)\right)$
\end{itemize}

从\ref{subsec-math-pre}节开始，我们将先介绍泛函分析的一些基本数学工具。接着利用这些数学工具研究：
\begin{itemize}
    \item 如何求解贝尔曼方程
    \item 函数与策略函数的相关性质
\end{itemize}

为了理解使用数学工具的必要性，在此处先抛出一个引例。事实上，如果能够解决这个问题，值函数的性质自然得到证明。


% 1.3.2.3 引例
\subsubsection{引利：连续逼近法}

目前为止，我们虽然推出了最优条件，但对于最为关注的初始目标——值函数$v$的形式（存在性与唯一性），还没有任何认识。上述递归形式问题的古典方法是
连续逼近法（method of successive approximations）。对于我们建立的贝尔曼方程：
$$
v\left(k\right)=\max _{k^{\prime}}\left[u\left(\gamma\left(k\right)-k^{\prime}\right)+\beta v\left(k^{\prime}\right)\right]
$$
我们的目标从这个泛函方程中求解值函数$v$。假设有一个备选结果$v_0\left(k\right)$符合贝尔曼方程，接着，定义一个新的值函数：
$$
v_1\left(k\right) = \max _{k^{\prime}}\left[u\left(\gamma\left(k\right)-k^{\prime}\right)+\beta v_0\left(k^{\prime}\right)\right]
$$
如果对所有$k\geq 0$都有$v_0\left(k\right)=v_1\left(k\right)$，那么我们幸运的解决了问题。但通常，$v_0\left(k\right) \neq v_1\left(k\right)$，继续迭代，
$$
v_{n+1}\left(k\right) = \max _{k^{\prime}}\left[u\left(\gamma\left(k\right)-k^{\prime}\right)+\beta v_n\left(k^{\prime}\right)\right],
    \quad n = 0,1,\dots
$$
我们感兴趣的是，当\textbf{迭代次数}$n$区域无穷时，$v_n$如何变化？

为了分析这一问题，首先，假设有一个可行（feasibale）策略$g_0\left(k\right)$，这个策略未必是最优的，因此被称为次优策略（Suboptimal Policy），
引入的目的是作为逼近方法的初始猜测。
定义$v_0\left(k\right)$是这个次优策略下的值：
$$
v_0\left(k_0\right) = \sum_{t=0}^{\infty} \beta^t u\left(\gamma\left(k_t\right)-g_0\left(k_t\right)\right)
$$
其中资本路径由$k_{t+1}=g_0\left(k_t\right)$生成。显然，有如下递归形式成立：
$$
v_0\left(k\right) = u\left(\gamma\left(k\right)-g_0\left(k\right)\right) + \beta v_0\left(g_0\left(k\right)\right)
$$
此时，我们通过假定一个可行的次优策略$g_0\left(k\right)$，得到了一个备选值函数$v_0\left(k\right)$。
接着，根据上面$n$的迭代法则，可得到：
$$
\begin{aligned}
    v_1\left(k\right) &= \max _{k^{\prime}}\left[u\left(\gamma\left(k\right)-k^{\prime}\right) + \beta v_0\left(k^{\prime}\right)\right] \\
                &\geq \left[u\left(\gamma\left(k\right)-g_0\left(k\right)\right) + \beta v_0\left(g_0\left(k\right)\right)\right] \\
                &= v_0\left(k\right)   
\end{aligned}
$$
其中大于等于号是因为，如果我们今天能取到方程最大化的最优值，它一定优于任意取的可行次优策略（计划者不会比始终遵循策略$g_0(k)$做得更差）。
类似地，继续进行迭代：
$$
\begin{aligned}
    v_2\left(k\right) &= \max _{k^{\prime}}\left[u\left(\gamma\left(k\right)-k^{\prime}\right) + \beta v_1\left(k^{\prime}\right)\right] \\
                &\geq \max _{k^{\prime}}\left[u\left(\gamma\left(k\right)-k^{\prime}\right) + \beta v_0\left(k^{\prime}\right)\right] \\
                &= v_1\left(k\right)   
\end{aligned}
$$
最终不难得到：
$$v_{n+1}\left(k\right) \geq v_{n}\left(k\right), n=0,1,\dots$$
需要注意,这里的$n$是迭代次数而不是时间标记。

回到贝尔曼方程上：
$$
v\left(k\right)=\max _{k^{\prime}}\left[u\left(\gamma\left(k\right)-k^{\prime}\right)+\beta v\left(k^{\prime}\right)\right]
$$
\textbf{定义其右侧为}：
\begin{equation}    \label{eq-operator}
    Tv\left(k\right) \equiv \max _{k^{\prime}} \left[u\left(\gamma\left(k\right)-k^{\prime}\right)+\beta v\left(k^{\prime}\right)\right]
\end{equation}
$T$是一个算子（Operator）：其输入是一个函数$v$，返回的是一个新的函数$Tv$，这就是\textbf{贝尔曼算子}（Bellman Operator）。
类似地，第$n+1$次迭代过程的\textbf{右侧}可以定义为：
$$
Tv_n\left(k\right) \equiv \max _{k^{\prime}} \left[u\left(\gamma\left(k\right)-k^{\prime}\right)+\beta v_n\left(k^{\prime}\right)\right]
$$
利用这一记号可得：$v_{n+1}=Tv_n, ~ n=0,1,\dots$。进一步将贝尔曼方程\textbf{左侧}纳入观察，
不难看出，对于贝尔曼方程的求解等价于求解如下的\textbf{不动点问题}（Fixed Point Problem）：
$$v = Tv$$

在次优策略下的迭代过程中我们已经得到：
$$v_{n+1}=Tv_n \geq v_n$$
与不动点问题相结合来看，我们现在的目标就是证明：当$n$趋于无穷大时，单调递增的函数序列$\{v_n\}$收敛，且极限$v$满足贝尔曼方程。
这样我们不仅证明了贝尔曼方程的可解性，而且通过迭代实现了一个算法（构造了一个不动点），得到数值解。为了研究类似于\ref{eq-operator}定义的算子，
我们需要运用相关数学工具。例如，为了证明算子$T$将一个函数空间$C$映射到其自身，我们必须选择合适的函数空间进行研究。
此外，我们可能还需要相应空间上算子$T$的不动点定理。\textbf{后续将看到，我们所需的两个核心工具就是压缩映射定理和Blackwell充分条件}。


% 1.3.3 值函数性质：数学准备
\subsection{值函数性质：数学准备} \label{subsec-math-pre}

上面的例子其实引出了两个问题：
\begin{itemize}
    \item 通过定义贝尔曼算子，我们将递归问题变为不动点问题$v=Tv$的求解，但是这个解在理论上是否存在和唯一？
    \item 给定初始猜测的$v_0\left(k\right)$，由算子$T$生成的值函数序列$\{v_n\}$是否收敛到$v$？即这一算法是否可行？
\end{itemize}

证明这两个问题需要做一些泛函分析的数学准备,例如，如何定义函数间的距离、函数序列的收敛性等。

% 1.3.3.1
\subsubsection{度量空间中的收敛性}

\begin{definition}
    \textbf{度量空间}（Metric Space）是由一个非空集合$S$和某一度量或称为距离（distance function）$d: S \times S \rightarrow \mathbb{R}$所定义，
    其中度量$d$性质如下：$\forall x,y,z \in S$
    \begin{itemize}
        \item $d\left(x,y\right) \geq 0$，当且仅当$x=y$时取等
        \item $d\left(x,y\right) = d\left(y,x\right)$
        \item $d(x,y) \leq d\left(x,z\right)+d\left(z,y\right) $
    \end{itemize}
\end{definition}
在这里，集合$S$远不止于数集，也可能是函数集（例如我们感兴趣的值函数集）等。因此，度量空间的定义比普通欧氏距离（Euclidean Distance）更复杂。
但显然，实数集$S = \mathbb{R}$和绝对距离$d\left(x,y\right)=\left|x-y\right|$构成一个度量空间。
\begin{example}
    下面是一些常见的度量空间例子：
    \begin{itemize}
        \item 任意有限维空间$\mathbb{R}^n$和其上定义的距离：
            \begin{gather*}
                d\left(x,y\right) = \left(\sum_{i=1}^{n} \left|x_i-y_i\right|^p\right)^{1/p}, \quad 1 \leq p < \infty \\
                d\left(x,y\right) = \max _i \left(\left|x_i-y_i\right|\right), \quad p=\infty
            \end{gather*}

        \item 定义在$\left[a,b\right]$上的连续函数集（常记为$C\left[a,b\right]$）和其上定义的距离：     
            \begin{gather*}
                d\left(x,y\right)=\left(\int_{a}^{b} \left|x\left(t\right)-y\left(t\right)\right|^{p}dt\right)^{1/p}, \quad 1 \leq p < \infty \\
                d\left(x,y\right) = \max_{a\leq t \leq b} \left(\left|x\left(t\right)-y\left(t\right)\right|\right), \quad p = \infty
            \end{gather*}

    \end{itemize}
\end{example}

\begin{definition}
    \textbf{赋范向量空间}（Normed Vector Space）是由一个向量空间$S$和范数$\| \cdot \|: S \rightarrow \mathbb{R}$所定义，
    其中范数性质如下：$\forall x,y \in S$和$a \in \mathbb{R}$
    \begin{itemize}
        \item \( \|x\| = 0 \)当且仅当 \(x=\mathbf{0}\)
        \item $\|x\|=|a|\cdot \|x\|$
        \item $\|x+y\| \leq \|x\|+\|y\|$
    \end{itemize}
\end{definition}
不难看出，如果令距离$d\left(x,y\right)=\|x-y\|$，那么任意赋范向量空间都是度量空间。
\begin{example}
    下面是一些常见的赋范向量空间例子：
    \begin{itemize}
        \item 任意有限维空间$\mathbb{R}^n$和其上定义的范数：
            \begin{gather*}
                \|x\|=\left(\sum_{i=1}^{n}|x_i|^p\right)^{1/p}, \quad 1\leq p < \infty \\
                \|x\|=\max_{i}\left(|x_i|\right), \quad p=\infty
            \end{gather*}

        \item 定义在$\left[a,b\right]$上的连续函数集$C\left[a,b\right]$和其上定义的范数：     
            \begin{gather*}
                \|x\|=\left(\int_{a}^{b}|x\left(t\right)|^{p}dt\right)^{1/p}, \quad 1 \leq p < \infty \\
                \|x\|=\max _{a\leq t \leq b}\left(|x\left(t\right)|\right), \quad p = \infty
            \end{gather*}
    \end{itemize}
\end{example}

\begin{definition}[\textbf{收敛性}]
    我们称$S$中的序列$\{x_n\}_{n=0}^{\infty}$收敛到$x \in S$：如果对$\forall \epsilon > 0$，存在$N_{\epsilon}$，使得：
    $$d\left(x_n,x\right)<\epsilon, \quad n\geq N_{\epsilon}$$
\end{definition}
注意序列$\{x_n\}$不只限于实数（事实上在动态规划中，我们感兴趣的是函数序列）。但从另一角度看，
序列$\{x_n\}$收敛到$x \in S$等价于距离的实数序列$\{d\left(x_n,x\right)\}$收敛到0。

\begin{definition}[\textbf{柯西序列}]
    $S$中的序列$\{x_n\}_{n=0}^{\infty}$被称为柯西序列（Cauchy Sequence）如果其满足柯西准测：
    对$\forall \epsilon > 0$，存在$N_{\epsilon}$，使得：
    $$d\left(x_n,x_m\right) < \epsilon, \quad n,m \geq N_{\epsilon}$$
\end{definition}

\begin{theorem}
    所有收敛序列都是柯西序列
\end{theorem}
\begin{theorem}
    所有柯西序列都有界
\end{theorem}
定义柯西序列的好处在于，可以只利用序列本身信息进行判断，而不用知道实际极限。
但是需要注意，\textbf{并非所有柯西序列都收敛}。实际上，这与序列所处空间的\textbf{完备性}（Completeness）密切相关。
\begin{example}
    考虑$\left[0,1\right]$上的连续函数集合$C\left[0,1\right]$。定义$L^1$度量：
    $d\left(f,g\right) = \int_{0}^{1}\left|f\left(x\right)-g\left(x\right)\right|dx$。
    构造函数序列$\{f_n\}$：
    $$
    f_n(x)= \begin{cases}0 & x \in\left[0, \frac{1}{2}-\frac{1}{n}\right] \\ 
        n\left(x-\left(\frac{1}{2}-\frac{1}{n}\right)\right) & x \in\left(\frac{1}{2}-\frac{1}{n}, \frac{1}{2}\right] 
        \\ 1 & x \in\left(\frac{1}{2}, 1\right] \end{cases}
    $$
    由于对任意$m>n$，成立：
    $$
    d\left(f_m, f_n\right)=\int_0^1\left|f_m(x)-f_n(x)\right| d x \leq \frac{1}{n} \rightarrow 0 \quad(n \rightarrow \infty)
    $$
    $\{f_n\}$在$L^1$度量下是柯西序列，但其收敛到极限为：
    $$
    f(x)= \begin{cases}0 & x \in\left[0, \frac{1}{2}\right] 
        \\ 1 & x \in\left(\frac{1}{2}, 1\right]\end{cases}
    $$
    在$x=\frac{1}{2}$处不连续，因此极限函数$f \notin C[0,1]$。从而 $\left\{f_n\right\}$ 在 $C[0,1]$ 中无极限，不是收敛序列。
\end{example}

\begin{definition}[\textbf{完备性}]
    如果$S$中任意柯西序列都收敛到$S$中某一极限，则称度量空间$\left(S,d\right)$是完备的（Complete）。完备的赋范向量空间被称为巴拿赫（Banach）空间。
\end{definition}

\begin{theorem} \label{Complete-CX}
    令$X \subseteq \mathbb{R}^n$，定义了上确界范数：$\|f\|=\sup _{x \in X}(|f(x)|)$的所有有界连续函数$f: X \rightarrow \mathbb{R}$的集合
    $C(X)$是一个完备空间。
\end{theorem}

定义完备空间的方便之处在于，只要其中的序列符合柯西准测，那么它在这一空间中就存在极限。


% 1.3.3.2 CMT
\subsubsection{压缩映射定理}

我们已经指出，贝尔曼方程\ref{eq-bellman}的本质是求解一个泛函中的不动点问题：$v=Tv$。其中$T$是贝尔曼算子：
$$
Tv\left(k\right) \equiv \max_{k^{\prime}} \left[u\left(\gamma\left(k\right)-k^{\prime}\right)+\beta v\left(k^{\prime}\right)\right]
$$
同时,从算法层面（寻找不动点）我们已经构造了递增的值函数序列$\{v_n\}$：
\begin{equation*}
    \begin{aligned}
        Tv_n\left(k\right) &\equiv \max _{k^{\prime}} \left[u\left(\gamma\left(k\right)-k^{\prime}\right)+\beta v_n\left(k^{\prime}\right)\right] \\
        v_{n+1} &= Tv_n \geq v_n
    \end{aligned}
\end{equation*}

我们希望证明，随着迭代次数$n$足够大，值函数序列$\{v_n\}$收敛到唯一极限$v$。这需要用到\textbf{压缩映射定理}（Contraction Mapping Theorem）。
什么是压缩映射？和前面定义的度量空间和完备性有什么关系？

\begin{definition}[\textbf{压缩映射}]
    令$\left(S,d\right)$为度量空间，并定义$T:S\rightarrow S$。如果对某一$\beta \in \left(0,1\right)$，成立：
    $$
    d\left(Tx,Ty\right) \leq \beta d\left(x,y\right)
    $$
    则称$T$为压缩映射，$\beta$为压缩系数。
\end{definition}
\begin{example}
    令$S=\left[a,b\right]$，$d\left(x,y\right)=\left|x-y\right|$。当对某一$\beta \in \left(0,1\right)$：
    $$\left|Tx-Ty\right|\leq \beta \left|x-y\right|$$
    则$T:S\rightarrow S$为压缩映射。
    考虑$a=0,b=1$，$Tx=0.2+0.6x$：
    $$\left|Tx-Ty\right|=0.6\left(x-y\right)\leq \beta\left|x-y\right|,\quad \forall \beta \in [0.6,1)$$
    因此$T$是压缩映射。
\end{example}

\begin{definition}
    若$x \in S$满足：$Tx=x$，则称$x$为不动点。
\end{definition}

\begin{theorem}[Contraction Mapping Theorem]   \label{th-cmt}
    令$\left(S,d\right)$为完备度量空间，令$T:S\rightarrow S$为压缩映射，压缩系数为$\beta$，则
    \begin{itemize}
        \item $T$在$S$中有唯一不动点；
        \item 对任意$x_0 \in S$，成立：
            $d\left(T^nx_0, x\right) \leq \beta^n d\left(x_0,x\right), \quad n=0,1,2,\dots$
    \end{itemize}
\end{theorem}

\begin{proof}
    先证明不动点存在，接着利用压缩映射定义。前者分三步证明：

    \textbf{Step1：证明迭代算子$T$后形成的序列是一个柯西序列}

    任取一个初始$x_0 \in S$，令$x_n=T^n x_0$，进而得到了序列$\{x_n\}_{n=0}^{\infty}$并且满足$x_{n+1}=Tx_n$。
    由于$T$是压缩映射，因此有：$d\left(x_2, x_1\right)=d\left(T x_1, T x_0\right) \leq \beta d\left(x_1, x_0\right)$,
    以此类推，对$n\geq 1$成立：
    $$
    d\left(x_{n+1}, x_n\right) \leq \beta^n d\left(x_1, x_0\right)
    $$
    因此对任意$m>n$，应用距离的三角不等式性质：
    $$
    \begin{aligned}
    d\left(x_m, x_n\right) & \leq d\left(x_m, x_{m-1}\right)+\cdots+d\left(x_{n+2}, x_{n+1}\right)+d\left(x_{n+1}, x_n\right) \\
    & \leq\left[\beta^{m-1}+\cdots+\beta^{n+1}+\beta^n\right] d\left(x_1, x_0\right) \\
    & =\beta^n\left[\beta^{m-n-1}+\cdots+\beta+1\right] d\left(x_1, x_0\right) \\
    & \leq \frac{\beta^n}{1-\beta} d\left(x_1, x_0\right)
    \end{aligned}
    $$
    由于$\beta \in \left(0,1\right)$，显然$\{x_n\}$为柯西序列。又由于$S$是完备的，因此$x_n\rightarrow x \in S$

    \textbf{Step2：证明极限是$T$的不动点}

    再次运用不三角不等式和压缩映射的定义：
    $$
    \begin{aligned}
        d(T x, x) & \leq d\left(T x, T^n x_0\right)+d\left(T^n x_0, x\right) \\
        &= d\left(Tx,T\left(T^{n-1}x_0\right)\right) + \left(T^n x_0, x\right) \\
        & \leq \beta d\left(x, T^{n-1} x_0\right)+d\left(T^n x_0, x\right)
    \end{aligned}
    $$
    显然当$n\rightarrow \infty$，右边两项的距离都趋于0，因此：
    $$d\left(Tx,x\right)=0, \quad Tx = x$$
    从而序列$\{x_n\}$的极限$x$是$T$的不动点。

    \textbf{Step3：证明不动点的唯一性}

    反证，假设存在另一不动点$Tx^{\prime}=x^{\prime},x^{\prime} \neq x$，则：
    $$
    0<\delta \equiv d\left(x,x^{\prime}\right)=d\left(Tx,Tx^{\prime}\right)\leq \beta d\left(x,x^{\prime}\right)=\beta \delta
    $$
    但由于$\beta \in \left(0,1\right)$，显然矛盾，从而唯一性得证。

    最后，由压缩映射的定义立得：
    $$
    d\left(T^nx_0,x\right) = d\left(T\left(T^{n-1}x_0\right),Tx\right) \leq \beta d\left(T^{n-1}x_0,x\right)
    $$
\end{proof}

压缩映射定理\ref{th-cmt}表明：
\begin{itemize}
    \item \textbf{完备空间}上定义的压缩映射存在唯一不动点
    \item 从\textbf{完备空间}$S$中任意初始点$x_0 \in S$出发，在其上迭代压缩映射$T$，可以无限接近唯一的不动点，即我们找到了一个求解不动点的算法
\end{itemize}

再次回顾目标：我们通过迭代贝尔曼算子：
$$
Tv\left(k\right) \equiv \max_{k^{\prime}} \left[u\left(\gamma\left(k\right)-k^{\prime}\right)+\beta v\left(k^{\prime}\right)\right]
$$
得到了递增的值函数序列$v_{n+1}=Tv_n \geq v_n$，我们希望这一序列收敛到贝尔曼方程的不动点$v=Tv$。我们已经指出，首先需要证明贝尔曼方程（不动点问题）确实存在唯一解；
其次去证明递增的值函数序列$\{v_n\}$收敛到不动点。显然，\textbf{如果能够证明贝尔曼算子是压缩映射}，那么直接应用压缩映射定理即可。
但是难点在于利用压缩映射的定义证明贝尔曼算子是压缩映射的难度较大，在这里给出一个充分条件：
\begin{theorem}[Blackwell充分条件]  \label{th-blackwell}
    令$X \subseteq \mathbb{R}^n$，令$B(X)$为有界函数集合$f: X \rightarrow \mathbb{R}$，其上定义了上确界范数（sup norm）：$\|f\|=\sup _{x \in X}(|f(x)|)$。
    令$T:B\left(X\right)\rightarrow B\left(X\right)$，且满足：
    \begin{itemize}
        \item 单调性：对$\forall f, g \in B(X)$，$f \leq g$， 则$T f \leq T g$
        \item discounting：存在$\beta \in \left(0,1\right)$，使得：$T(f+a) \leq T f+\beta a, \quad \forall f \in B(X)$。其中，
            函数$\left(f+a\right)\equiv f\left(x\right)+a$，$a \geq 0$。
    \end{itemize}
    那么，$T$为压缩映射。
\end{theorem}

\begin{proof}
    令$v, w \in B(X)$，从而：
    $$
    \begin{aligned}
        v(x) & =v(x)+w(x)-w(x) \\
        & \leq w(x)+|v(x)-w(x)| \\
        & \leq w(x)+\|v-w\|, \quad \text { for all } x \in X
    \end{aligned}
    $$
    由于映射$T$满足单调性和discounting，因此：
    $$
    T v \leq T(w+\|v-w\|) \leq T w+\beta\|v-w\|
    $$
    同理，
    $$
    T w \leq T(v+\|v-w\|) \leq T v+\beta\|v-w\|
    $$
    根据这两个不等式可知：
    $$
    |T v(x)-T w(x)| \leq \beta\|v-w\|, \quad \text { for all } x \in X
    $$
    也即，
    $$
    \|T v-T w\| \leq \beta\|v-w\|
    $$
    从而$T$是压缩映射。
\end{proof}

% 1.3.3.3
\subsubsection{应用CMT}
重新考虑最优增长模型。我们定义了贝尔曼方程右边为贝尔曼算子：
$$
Tv\left(k\right) \equiv \max_{k^{\prime}} \left[u\left(\gamma\left(k\right)-k^{\prime}\right)+\beta v\left(k^{\prime}\right)\right]
$$
\begin{itemize}
    \item $T$满足单调性：若$v \leq w$，则
        $$u\left(\gamma\left(k\right)-k^{\prime}\right)+\beta v\left(k^{\prime}\right) \leq 
            u\left(\gamma\left(k\right)-k^{\prime}\right)+\beta w\left(k^{\prime}\right), \quad \forall x
        $$
        从而
        $$
        \max _{k^{\prime}}\left[u\left(\gamma\left(k\right)-k^{\prime}\right)+\beta v\left(k^{\prime}\right)\right] \leq 
        \max _{k^{\prime}}\left[u\left(\gamma\left(k\right)-k^{\prime}\right)+\beta w\left(k^{\prime}\right)\right]
        $$
        因此
        $$Tv\left(k\right)\leq Tw\left(k\right)$$
    \item T满足discounting：对$\forall a \geq 0$，成立：
        $$
        \begin{aligned}
            T(v+a)(k) & =\max _x[u(f(k)-x)+\beta(v(x)+a)] \\
            & =\max _x[u(f(k)-x)+\beta v(x)]+\beta a \\
            & =T v(k)+\beta a
        \end{aligned}
        $$
\end{itemize}
因此，贝尔曼算子$T$满足Blackwell充分条件，从而是压缩映射。

再次回忆此前提到所面临的两个问题：
\begin{itemize}
    \item 理论层面，贝尔曼方程$v=Tv$有解吗，即是否存在唯一不动点
    \item 算法层面，构造的值函数序列是否收敛
\end{itemize}

在证明了贝尔曼算子$T$是压缩映射后，直接应用压缩映射定理，似乎我们已经完成了证明。但实际上，还存在两个确实缺陷：
一个存在于证明贝尔曼算子压缩性质（即运用Blackwell充分条件）的过程中；另一个则存在于运用压缩映射定理的过程中。具体看：

（1）\textbf{Blackwell充分条件应用的前提}：算子$T$是定义在\textbf{有界函数空间}上的。

对于递归形式下的最优增长模型，我们所分析的也是一个函数空间，
即$S$中的元素是值函数$v$。但是，值函数是否有界？答案是肯定的，这与我们对生产函数和效用函数的假设密切相关。

根据Inada条件、$\beta < 1$和消费与资本的非负性，由级数收敛性知识，下界的存在是显然的。
对于上界：由生产函数的凹性，资本积累存在上界，记为$k \leq k^{\star}$，从而$f\left(k\right) \leq f\left(k^{\star}\right)$；
消费的最大可能取值则为$f\left(k\right)$，即下期资本取0，从而有：
$u\left(c\right)\leq u\left(f\left(k\right)\right)\leq u\left(f\left(k^{\star}\right)\right)$。因此，
$$
v\left(k\right) = \sum_{t=0}^{\infty} \beta^t u\left(c\right) \leq 
    \sum_{t=0}^{\infty} \beta^t u\left(f\left(k^{\star}\right)\right)
    = \frac{u\left(f\left(k^{\star}\right)\right)}{1-\beta} < \infty
$$
从而值函数有界，可以对贝尔曼算子应用Blackwell充分条件。

（2）\textbf{压缩映射定理应用的前提}：压缩映射$T$是定义在\textbf{完备空间}$\left(S,d\right)$上的。

值函数是有界的，但是值函数空间否是完备的呢？我们曾在定理\ref{Complete-CX}中指出，
定义了上确界范数$\|f\|=\sup _{x \in X}(|f(x)|)$的\textbf{有界连续函数集合}空间$C\left(X\right)$是一个完备空间。
那么，只要能证明\textbf{值函数是连续的}，我们就能在Blackwell条件的基础上进一步应用压缩映射定理：
证明不动点的存在性和唯一性，并且运用其推广证明值函数序列的收敛性。


% 1.3.4 最大值原理与值函数性质
\subsection{最优性原理与值函数性质} \label{subsec-principal-opt}

值函数连续性和相关性质的证明相对较为复杂，我们通过引入最优性原理（Principal of Optimality）进而展开。

% 1.3.4.1最优性原理
\subsubsection{最优性原理}

首先，我们可以写出更一般化形式的序列问题（Sequential Problem，$\mathbf{SP}$）：
$$
\sup _{\left\{x_{t+1}\right\}_{t=0}^{\infty}} \sum_{t=0}^{\infty} \beta^t F\left(x_t, x_{t+1}\right)
$$
满足：
$$
\begin{aligned}
    x_{t+1} & \in \Gamma\left(x_t\right), \quad \forall t\\
    x_0 & \in X \quad \text { 给定 }
\end{aligned}
$$
其中，$X$是$x$所有可能取值集合；$\Gamma : X \rightarrow X$是描述可行性的对应（Correspondence），即
给定当前状态$x_t$，下期$x_{t+1}$的所有可行取值；$F$是当期回报函数。我们称从$x_0$开始，任意可行的序列$\{x_{t+1}\}_{t=0}^{\infty}$为一个可行计划（Feasible Plans）。
记所有可行计划构成的集合为：
$$\pi\left(x_0\right)=\left\{\left\{x_{t+1}\right\}_{t=0}^{\infty} \quad: \quad x_{t+1} \in \Gamma\left(x_t\right), \quad t=0,1,2, \ldots\right\}$$
其元素是一个序列，记为$\boldsymbol{x}\in\pi\left(x_0\right)$。

其次，我们写出与$SP$相对应的递归问题（$R$），或称为泛函方程（Functional Equation，$\mathbf{FE}$）：
$$
v(x)=\sup _{y \in \Gamma(x)}[F(x, y)+\beta v(y)] \quad \forall x \in X
$$
$x$是状态变量，$y$是控制变量。

$SP$与$FE$是否等价？更一般的想法是：
\begin{itemize}
    \item 问题$R$的解（函数）$v$在$x=x_0$处的取值达到$SP$的上界
    \item 序列$\{x_{t+1}\}_{t=0}^{\infty}$达到问题$SP$上界，当且仅当
        $v\left(x_t\right)=F\left(x_t, x_{t+1}\right)+\beta v\left(x_{t+1}\right), \quad t=0,1,2, \ldots$
\end{itemize}

贝尔曼将这一想法称为\textbf{最优性原理}。当然，这一结果的成立需要一些假设。
\begin{assumption} \label{as-nonempty-plan}
    集合$\Gamma\left(x\right)$对任意$x\in X$非空
\end{assumption}
\begin{assumption} \label{as-lim-exist}
    对任意$\left\{x_0, x_1, \ldots\right\} \in \Pi\left(x_0\right), x_0 \in X$，
        目标极限$\lim\limits_{n\rightarrow \infty}\sum_{t=0}^{n}\beta^t F\left(x_t,x_{t+1}\right)$存在（可能为无穷）
\end{assumption}
不难看出，假设\ref{as-lim-exist}成立的一个充分条件是：$F\left(x,y\right)$有界（bounded）且$0<\beta<1$。
对任意$n=0,1,\dots$，定义$u_n: \pi\left(x_0\right)\rightarrow \mathbb{R}$：
$$
u_n(\boldsymbol{x}) \equiv \sum_{t=0}^n \beta^t F\left(x_t, x_{t+1}\right), \quad \boldsymbol{x}=\{x_0,x_1,\dots\} \in \pi\left(x_0\right)
$$
由假设\ref{as-lim-exist}，可定义极限：
\begin{equation*}
    u\left(\boldsymbol{x}\right) \equiv \lim _{n\rightarrow\infty}u_n\left(\boldsymbol{x}\right)
\end{equation*}
如果假设\ref{as-nonempty-plan}和假设\ref{as-lim-exist}均成立，那么序列问题$SP$是良定义的。因此我们可以定义\textbf{上确界函数}：
\begin{equation*}
    v^*\left(x_0\right)\equiv \sup _{\boldsymbol{x}\in\pi\left(x_0\right)} u\left(\boldsymbol{x}\right)
\end{equation*}
这样，$v^*\left(x_0\right)$是$SP$问题的上确界，即符合上确界定义：
\begin{equation*}
    \begin{aligned}
        &\forall \epsilon > 0, \forall \boldsymbol{x}\in \pi\left(x_0\right): \quad v^*\left(x_0\right)\geq u\left(\boldsymbol{x}\right) \\
        &\forall \epsilon > 0, \exists \boldsymbol{x}\in\pi\left(x_0\right): \quad v^*\left(x_0\right)\leq u\left(\boldsymbol{x}\right)+\epsilon
    \end{aligned}
\end{equation*}

我们希望研究$SP$问题的上确界\textbf{函数}$v^*$和$FE$问题的解$v$的关系。需要注意：$v^*$是唯一定义的，但$FE$问题的解$v$可能有很多个。
简单起见，以下不考虑无穷大的情形，即考虑$\left|v^*\left(x_0\right)\right| < \infty$。那么，$v^*$是$FE$问题的解等价于如下条件成立：
$$
\begin{aligned}
    &\forall y \in \Gamma\left(x_0\right),\quad v^*\left(x_0\right)\geq F\left(x_0,y\right) + \beta v^*\left(y\right)\\
    &\exists y\in \Gamma\left(x_0\right),\quad v^*\left(x_0\right)\leq F\left(x_0,y\right) + \beta v^*\left(y\right), \quad \forall \epsilon > 0
\end{aligned}
$$
为了证明$v^*$满足$FE$，引入如下引理：
\begin{lemma}
    令$X, \Gamma, F, \beta$符合假设\ref{as-lim-exist}，
    则对$\forall x_0 \in X, \forall \left(x_0,x_1,\dots\right)=\boldsymbol{x}\in\pi\left(x_0\right)$，成立：
    \begin{equation*}
        u\left(\boldsymbol{x}\right)=
            F\left(x_0,x_1\right)+\beta u\left(\boldsymbol{x}^{\prime}\right), \quad \boldsymbol{x}^{\prime}=\left(x_1,x_2,\dots\right)
    \end{equation*}
\end{lemma}
\begin{theorem}[SP-FE] \label{th-solution-sp-fe}
    令$X, \Gamma, F, \beta$符合假设\ref{as-nonempty-plan}和假设\ref{as-lim-exist}，则上确界函数$v^*$满足$FE$。
\end{theorem}
\begin{proof}
    前面提到，上确界函数$v^*$是$FE$问题的解等价于成立如下两个条件：
    $$
    \begin{aligned}
        &\forall y \in \Gamma\left(x_0\right),\quad v^*\left(x_0\right)\geq F\left(x_0,y\right) + \beta v^*\left(y\right)\\
        &\exists y\in \Gamma\left(x_0\right),\quad v^*\left(x_0\right)\leq F\left(x_0,y\right) + \beta v^*\left(y\right), \quad \forall \epsilon > 0
    \end{aligned}
    $$
    
    先证明第一个条件：$\forall y \in \Gamma\left(x_0\right),\quad v^*\left(x_0\right)\geq F\left(x_0,y\right) + \beta v^*\left(y\right)$。

    由$v^*\left(x_0\right)\equiv \sup _{\boldsymbol{x}\in\pi\left(x_0\right)} u\left(\boldsymbol{x}\right)$，
    根据上确界的定义，成立如下两条：
    \begin{equation*}
        \begin{aligned}
            &\forall \epsilon > 0, \forall \boldsymbol{x}\in \pi\left(x_0\right): \quad v^*\left(x_0\right)\geq u\left(\boldsymbol{x}\right) \\
            &\forall \epsilon > 0, \exists \boldsymbol{x}\in\pi\left(x_0\right): \quad v^*\left(x_0\right)\leq u\left(\boldsymbol{x}\right)+\epsilon
        \end{aligned}
    \end{equation*}
    根据后者，任取$x_1 \in \Gamma\left(x_0\right)$，任取$\epsilon > 0$：
    $$\exists \boldsymbol{x}^{\prime}=\left(x_1,x_2,\dots\right) \in \pi\left(x_1\right), 
        \quad u\left(\boldsymbol{x}^{\prime}\right)\geq v^*\left(x_1\right)-\epsilon$$
    又由于$\boldsymbol{x}=\left(x_0,x_1,\dots\right)\in \pi\left(x_0\right)$，根据前者和上述引理：
    $$
    v^*\left(x_0\right)\geq u\left(\boldsymbol{x}\right)=F\left(x_0,x_1\right)+\beta u\left(\boldsymbol{x}^{\prime}\right)
        \geq F\left(x_0,x_1\right)+\beta v^*\left(x_1\right)-\beta \epsilon
    $$
    由于$\epsilon$任意，可知：
    $$
    \forall y \in \Gamma\left(x_0\right),\quad v^*\left(x_0\right)\geq F\left(x_0,y\right) + \beta v^*\left(y\right)
    $$

    再证明第二个条件：$\exists y\in \Gamma\left(x_0\right),\quad v^*\left(x_0\right)\leq F\left(x_0,y\right) + \beta v^*\left(y\right), \quad \forall \epsilon > 0$。

    任取$x_0\in X$和$\epsilon>0$，根据上确界函数定义的第二条和引理可知：
    $$
    \exists \boldsymbol{x}=\left(x_0,x_1,\dots\right)\in \pi\left(x_0\right), 
        \quad v^*\left(x_0\right)\leq u\left(\boldsymbol{x}\right)+\epsilon = F\left(x_0,x_1\right)+\beta u\left(\boldsymbol{x}^{\prime}\right)+\epsilon
    $$
    其中，$\boldsymbol{x}^{\prime}= \left(x_1,x_2,\dots\right)$。从而，结合上确界定义的第一条可知：
    $$
    v^*\left(x_0\right)\leq F\left(x_0,x_1\right)+\beta v^*\left(x_1\right)+\epsilon 
    $$
    由$\epsilon$任意性和$x_1\in\Gamma\left(x_0\right)$（因为$x_1$在计划中）：
    $$
    \exists y\in \Gamma\left(x_0\right),\quad v^*\left(x_0\right)\leq F\left(x_0,y\right) + \beta v^*\left(y\right), \quad \forall \epsilon > 0
    $$
    从而我们证明了$v^*$是$FE$问题的解。
\end{proof}

至此，我们证明了$SP$的解满足$FE$，但反过来是否成立？一般来说不一定，虽然$v^*$是$FE$的解，但$FE$可能存在很多解。
但下面的定理表明：若$FE$满足特定有界条件，那么$v=v^*$是其唯一解。

\begin{theorem}[Partial Converse: FE-SP] \label{th-solution-fe-sp}
    令$X, \Gamma, F, \beta$符合假设\ref{as-nonempty-plan}和假设\ref{as-lim-exist}。如果值函数$v$是$FE$的一个解且满足
    $\lim\limits_{n\rightarrow \infty}\beta^n v\left(x_n\right)=0, 
    \quad \forall x_0\in X, \forall \boldsymbol{x}=\left(x_0,x_1,\dots\right)\in \pi\left(x_0\right)$，
    则$v=v^*$。其中：
    \begin{gather*}
        v^*\left(x_0\right) \equiv \sup _{\boldsymbol{x}\in\pi\left(x_0\right)} u\left(\boldsymbol{x}\right) \\
        u\left(\boldsymbol{x}\right) \equiv \lim _{n\rightarrow\infty}u_n\left(\boldsymbol{x}\right) \\
        u_n(\boldsymbol{x}) \equiv \sum_{t=0}^n \beta^t F\left(x_t, x_{t+1}\right), \quad \boldsymbol{x}=\{x_0,x_1,\dots\} \in \pi\left(x_0\right)
    \end{gather*}
\end{theorem}
\begin{proof}
    方便起见仍考虑$v\left(x_0\right)$有限情形。我们此前提到此时$v$是$FE$的解等价于如下条件成立：
    $$
    \begin{aligned}
        &\forall y \in \Gamma\left(x_0\right),\quad v\left(x_0\right)\geq F\left(x_0,y\right) + \beta v\left(y\right)\\
        &\exists y\in \Gamma\left(x_0\right),\quad v\left(x_0\right)\leq F\left(x_0,y\right) + \beta v\left(y\right), \quad \forall \epsilon > 0
    \end{aligned}
    $$
    我们的目标是证明$FE$的这一解$v$是$SP$的解，即证明
    $v\left(x_0\right)=\sup _{\boldsymbol{x}\in\pi\left(x_0\right)} u\left(\boldsymbol{x}\right)$，也即往证$FE$的解$v$符合$SP$中上确界的定义：
    \begin{equation*}
        \begin{aligned}
            &\forall \epsilon > 0, \forall \boldsymbol{x}\in \pi\left(x_0\right): \quad v\left(x_0\right)\geq u\left(\boldsymbol{x}\right) \\
            &\forall \epsilon > 0, \exists \boldsymbol{x}\in\pi\left(x_0\right): \quad v\left(x_0\right)\leq u\left(\boldsymbol{x}\right)+\epsilon
        \end{aligned}
    \end{equation*}
    注意到对任意$\boldsymbol{x}\in \pi\left(x_0\right)$，成立：
    \begin{equation*}
        \begin{aligned}
            v\left(x_0\right) &\geq F\left(x_0,x_1\right) + \beta v\left(x_1\right) \\
                &\geq  F\left(x_0,x_1\right) + \beta F\left(x_1,x_2\right) + \beta^2v\left(x_2\right) \\
                &\geq \cdots \\
                &\geq u_n\left(\boldsymbol{x}\right)+\beta^{n+1}v\left(x_{n+1}\right)
        \end{aligned}
    \end{equation*}
    取极限并利用
    $\lim\limits_{n\rightarrow \infty}\beta^n v\left(x_n\right)=0, 
        \quad \forall x_0\in X, \forall \boldsymbol{x}=\left(x_0,x_1,\dots\right)\in \pi\left(x_0\right)$
    可知上确界定义的前一半成立：
    \begin{equation*}
        \forall \epsilon > 0, \forall \boldsymbol{x}\in \pi\left(x_0\right): 
            \quad v\left(x_0\right)\geq u\left(\boldsymbol{x}\right)
    \end{equation*}
    下面证明上确界定义的后一半，即：
    \begin{equation*}
        \forall \epsilon > 0, \exists \boldsymbol{x}\in\pi\left(x_0\right): \quad v\left(x_0\right)\leq u\left(\boldsymbol{x}\right)+\epsilon  
    \end{equation*}
    对此，任取一个$\epsilon>0$，从正实数集中选择$\{\delta_t\}_{t=1}^{\infty}$，使得$\sum_{t=1}^{\infty}\beta^{t-1}\delta_t\leq \epsilon /2$。
    由与$v$是$FE$的解，即：
    \begin{equation*}
        \exists y\in \Gamma\left(x_0\right),\quad v\left(x_0\right)\leq F\left(x_0,y\right) + \beta v\left(y\right), \quad \forall \epsilon > 0
    \end{equation*}
    从而可以取到$x_1\in\Gamma\left(x_0\right), x_2\in\Gamma\left(x_1\right),\cdots$，使：
    \begin{equation*}
        v\left(x_t\right)\leq F\left(x_t,x_{t+1}\right)+\beta v\left(x_{t+1}\right)+\delta_{t+1},\quad t=0,1,\cdots
    \end{equation*}
    也就是说，$\exists\boldsymbol{x}=\left(x_0,x_1,\cdots\right)\in \pi\left(x_0\right)$，使：
    \begin{equation*}
        \begin{aligned}
            v\left(x_0\right) &\leq \sum_{t=0}^{n}\beta^t F\left(x_t,x_{t+1}\right) + 
                \beta^{n+1}v\left(x_{n+1}\right) + \left(\delta_1+\cdots+\beta^n \delta_{n+1}\right)   \\
                &\leq u_n\left(\boldsymbol{x}\right)+\beta^{n+1}v\left(x_{n+1}\right)+\epsilon/2, \quad n=1,2,\cdots
        \end{aligned}
    \end{equation*}
    当$n$足够大，由条件$\lim_{n\rightarrow \infty}\beta^n v\left(x_n\right)=0$可知：
       $v\left(x_0\right)\leq u_n\left(\boldsymbol{x}\right)+\epsilon$。由于$\epsilon$的可知任意性，上确界后半的定义成立：
    \begin{equation*}
        \forall \epsilon > 0, \exists \boldsymbol{x}\in\pi\left(x_0\right): \quad v^*\left(x_0\right)\leq u\left(\boldsymbol{x}\right)+\epsilon
    \end{equation*}
    从而：
    \begin{equation*}
        v\left(x_0\right)=\sup _{\boldsymbol{x}\in\pi\left(x_0\right)} u\left(\boldsymbol{x}\right)
    \end{equation*}
    也即$FE$的解$v$是$SP$的解。
\end{proof}

至此，我们可以对序列问题（$SP$）和递归（泛函）问题（$FE$）解的关系进行一个总结。方便起见，再次写下这两个一般化的问题：
\begin{itemize}
    \item $SP$：
    $$\sup _{\left\{x_{t+1}\right\}_{t=0}^{\infty}} \sum_{t=0}^{\infty} \beta^t F\left(x_t, x_{t+1}\right),
    \quad x_{t+1}  \in \Gamma\left(x_t\right), \forall t, \quad x_0 \in X \quad \text { 给定 }
    $$
    \item $FE$：
    $$
    v(x)=\sup _{y \in \Gamma(x)}[F(x, y)+\beta v(y)] \quad \forall x \in X
    $$
\end{itemize}

在假设\ref{as-nonempty-plan}和假设\ref{as-lim-exist}下，我们通过构造极限函数的上确界找到了$SP$的解：
$$
v^*\left(x_0\right)=\sup _{\boldsymbol{x}\in\pi\left(x_0\right)}u\left(\boldsymbol{x}\right), 
u\left(\boldsymbol{x}\right)=\lim_{n\rightarrow\infty}\sum_{t=0}^{n}\beta^t F\left(x_t,x_{t+1}\right)
$$
同时证明了它至少也是$FE$的一个解，即$SP\rightarrow FE$。

另一方面，$FE$可能存在很多解$v$，我们证明了其中符合有界条件
$
\lim\limits_{n\rightarrow \infty}\beta^n v\left(x_n\right)=0, \quad \forall x_0\in X, \forall \boldsymbol{x}=\left(x_0,x_1,\dots\right)\in \pi\left(x_0\right)
$
的解正是$SP$的解$v^*$，即部分可逆。而$FE$的其余解一定违反了这一有界条件.


% 1.3.4.2 最优策略的关系
\subsubsection{最优策略间的关系}

最优性原理（定理\ref{th-solution-sp-fe}和定理\ref{th-solution-fe-sp}）给出了序列问题和递归（泛函）问题目标值函数间的关系。
我们希望进一步研究两个问题最优策略（Plan）间的关系。

\begin{definition}[最优计划]
    如果可行计划$\boldsymbol{x}^* = \{x_{t+1}^*\} \in\pi\left(x_0\right)$能达到$SP$问题的上确界，即
    $u\left(\boldsymbol{x}^*\right) \equiv \sum_{t=0}^{\infty}\beta^t F\left(x_t^*,x_{t+1}^*\right) = v^*\left(x_0\right)$，
    则称$\boldsymbol{x}^*$为从$x_0$开始的最优计划（Optimal Plan）。
\end{definition}
下面两个定理给出了最优计划和$FE$问题中满足$v=v^*$的策略方程间的关系。

\begin{theorem} \label{th-policy-sp-fe}
    令$X, \Gamma, F, \beta$符合假设\ref{as-nonempty-plan}和假设\ref{as-lim-exist}。令$\boldsymbol{x}^*\in \pi\left(x_0\right)$
    为一个达到$SP$问题上确界的可行计划，则:
    \begin{equation} \label{eq-opt-plan-sp}
        v^*\left(x_t^*\right) = F\left(x_t^*,x_{t+1}^*\right) + \beta v^*\left(x_{t+1}^*\right), \quad t=0,1,2,\cdots
    \end{equation}
\end{theorem}

\begin{proof}
    由于$\boldsymbol{x}^*$达到了上确界，因此由定理\ref{th-solution-sp-fe}，$v^*\left(x_t^*\right)$满足$FE$，即：
    \begin{equation*}
        \begin{aligned}
            v^*\left(x_0^*\right) &= u\left(\boldsymbol{x}^*\right) = F\left(x_0,x_1^*\right) 
                + \beta u\left(\boldsymbol{x^{*\prime}}\right) \\
            &\geq u\left(\boldsymbol{x}\right) = F\left(x_0,x_1\right)+\beta u\left(\boldsymbol{x}^{\prime}\right), \quad 
                \forall \boldsymbol{x}\in \pi\left(x_0\right)
        \end{aligned}
    \end{equation*}
    又由于$\left(x_1^*,x_2,x_3,\cdots\right)\in \pi\left(x_1^*\right)$，
    从而$\left(x_0,x_1^*,x_2,x_3,\cdots\right)\in \pi\left(x_0\right)$,
    进而有：
    $$u\left(\boldsymbol{x}^{*\prime}\right) \geq u\left(\boldsymbol{x}^{\prime}\right), \quad \forall \boldsymbol{x}\in \pi\left(x_1^*\right)$$
    从而$u\left(\boldsymbol{x}^{*\prime}\right)=v\left(x_1^*\right)$，带入上式就得到了$t=0$时定理的结论。其余可依此类推即可。
\end{proof}

定理\ref{th-policy-sp-fe}表明，$SP$问题的最优计划是$FE$的一个解。
下面一个定理则提供了相反方向的直觉：它表明，在一定的有界条件下，符合式\ref{eq-opt-plan-sp}的序列是一个最优计划。

\begin{theorem} \label{th-policy-fe-sp}
    令$X, \Gamma, F, \beta$符合假设\ref{as-nonempty-plan}和假设\ref{as-lim-exist}。
    令$\boldsymbol{x}^*\in \pi\left(x_0\right)$，且满足
    $$v^*\left(x_t^*\right) = F\left(x_t^*,x_{t+1}^*\right) + \beta v^*\left(x_{t+1}^*\right), \quad t=0,1,2,\cdots$$
    并且$\lim\limits_{t\rightarrow\infty}\sup\beta^t v^*\left(x_t^*\right)\leq 0$。那么，$\boldsymbol{x}^*$达到$SP$的上确界，即是最优计划。
\end{theorem}
\begin{proof}
    由条件：
    \begin{equation*}
        \begin{aligned}
            v^*\left(x_0\right) = u\left(\boldsymbol{x}^*\right)&= F\left(x_0^*,x_1^*\right)+\beta v^*\left(x_1^*\right) \\
                &= \cdots \\
                &= u_n\left(\boldsymbol{x}^*\right) + \beta^{n+1}v^*\left(x_{n+1}^*\right)
        \end{aligned}
    \end{equation*}
    又由有界条件可知：$v^*\left(x_0\right)\leq u\left(\boldsymbol{x}^*\right)$。
    由于$\boldsymbol{x}^*\in \pi\left(x_0\right)$：$v^*\left(x_0\right)\geq u\left(\boldsymbol{x}^*\right)$
    从而：$v^*\left(x_0\right) = u\left(\boldsymbol{x}^*\right)$，由定义得证。
\end{proof}

\begin{definition}
    如果对任意非空对应$G: X\rightarrow X$，成立$G\left(x\right)\subseteq \Gamma\left(X\right)$，则称$G$为策略对应（Policy Correspondence）。
    若$G$为单值的，我们称其为策略函数（Policy Function），记为$g$。如存在序列$\boldsymbol{x}=\left(x_0,x_1,\cdots\right)$满足
    $x_{t+1}\in G\left(x_t\right), t=0,1,2,\cdots$，则称$\boldsymbol{x}$由$G$从$x_0$生成。可以定义最优策略对应（Optimal Policy Correspondence）$G^*$：
    $$
    G^*(x)=\left\{y \in \Gamma(x): v^*(x)=F(x, y)+\beta v^*(y)\right\}
    $$
\end{definition}
定理\ref{th-policy-sp-fe}表明：每个最优计划$\{x_t^*\}$都由$G^*$生成；
定理\ref{th-policy-fe-sp}表明：每个由$G^*$生成的计划$\{x_t^*\}$，若满足有界条件，就是一个最优计划。


% 1.3.4.3 FE解的存在性
\subsubsection{最大化定理和解的存在与唯一性}

我们在\ref{subsec-math-pre}节最后提出，对于递归形式的最有增长问题，我们希望利用Blackwell定理证明$T$是压缩映射，进而运用压缩映射定理证明贝尔曼方程
$v=Tv$这一不动点问题解的存在和唯一性。我们已经证明了贝尔曼算子$T$是压缩映射，但是在运用压缩映射定理时遇到了困难：我们希望值函数处于连续空间，
确保空间的完备性，进而才能使用压缩映射定理。

在\ref{subsec-principal-opt}节，我们将动态规划写成了更一般化的形式，只要完成一般化形式下的证明，
那么新古典增长模型中相关结论自然成立。我们还需要一些额外假设和定理的辅助。方便起见，在此处写下我们所需研究的一般化形式泛函方程：
\begin{equation}    \label{eq-general-fe}
    v\left(x\right)=\max_{y\in\Gamma\left(x\right)}\left[F\left(x,y\right)+\beta v\left(y\right)\right]
\end{equation}

\begin{assumption}[凸性假设] \label{as-convex-X}
    状态变量所有可能取值集合$X$是$\mathbb{R}^l$的凸子集，可行性约束对应$\Gamma :X\rightarrow X$是非空、连续、紧急。
\end{assumption}

\begin{assumption}[有界回报]  \label{as-bounded-returns}
    函数$F: A\rightarrow \mathbb{R}$是有界连续的，且$0<\beta<1$。
\end{assumption}

显然，假设\ref{as-convex-X}和假设\ref{as-bounded-returns}蕴含了假设\ref{as-nonempty-plan}和假设\ref{as-lim-exist}。
在\ref{subsec-math-pre}节最后所提到，为了应用压缩映射定理，我们不仅希望值函数有界，更希望在完备空间中研究值函数，因此
一个自然的想法就是定义了上确界范数$\|f\|=\sup _{x\in X}|f\left(x\right)|$的有界连续函数$f: X\rightarrow \mathbb{R}$构成空间$C\left(X\right)$。
定义空间$C\left(X\right)$上的算子$T$：
\begin{equation}
    Tf\left(x\right)=\max_{y\in\Gamma\left(x\right)}\left[F\left(x,y\right)+\beta f\left(y\right)\right]
\end{equation}
由此，我们将式\ref{eq-general-fe}简化为不动点问题$v=Tv$。
我们将借助最大化定理（Maximum Theorem）（定理\ref{th-max-th}）、Blackwell充分条件（定理\ref{th-blackwell}）和
压缩映射定理（定理\ref{th-cmt}）证明：算子$T$将$C\left(X\right)$中的元素（函数）映射到$C\left(X\right)$上，
即$T: C\left(X\right)\rightarrow C\left(X\right)$，从而$T$在有界连续函数空间$C\left(X\right)$上存在唯一不动点。

\begin{theorem}[Maximum Theorem]    \label{th-max-th}
    令$X\subseteq\mathbb{R}^l$，$Y\subseteq\mathbb{R}^m$，令$f: X\times Y \rightarrow \mathbb{R}$是连续函数，
    令$\Gamma: X\rightarrow Y$是紧值连续对应（Correspondence），那么：
    \begin{itemize}
        \item 从$X\rightarrow R$的函数$h\left(x\right)\equiv \max _{y\in\Gamma\left(x\right)}f\left(x,y\right)$是连续的
        \item 由$G(x) \equiv \underset{y \in \Gamma(x)}{\operatorname{argmax}} f(x, y)=\{y \in \Gamma(x): f(x, y)=h(x)\}$
        给定的解对应$G: X\rightarrow Y$是非空、紧值和上半连续的。
    \end{itemize}
\end{theorem}
对于形如：
\begin{equation*}
    h\left(x\right)\equiv \max_{y\in\Gamma\left(x\right)}f\left(x,y\right)
\end{equation*}
最大化定理给出了$h\left(x\right)$和相应最优$y$值的集合$G\left(x\right)$随$x$连续变化的条件。
同时，若$f\left(x,y\right)$关于$y$严格凹，并且$\Gamma\left(x\right)$是凸的，那么$G$是单值的，可记为函数$g\left(x\right)$。

\begin{theorem}[$FE$问题解的存在及唯一性] \label{th-fe-solution-exist}
    令$X, \Gamma, F, \beta$满足假设\ref{as-convex-X}和假设\ref{as-bounded-returns}，
    $C\left(X\right)$为定义了上确界范数的有界连续函数空间，那么：
    \begin{itemize}
        \item 算子$T$将$C\left(X\right)$映射到自身
        \item $T$有唯一不动点$v\in C\left(X\right)$
        \item 任意给定$v_0\in C\left(X\right)$，成立：$\|T^nv_0-v\|\leq \beta^n \|v_0-v\|$
    \end{itemize}
\end{theorem}

\begin{proof}
    \textbf{Step1：证明算子$T$将有界连续函数空间映射到自身}

    注意到在假设\ref{as-convex-X}和假设\ref{as-bounded-returns}下，$F$和$f$均有界连续，集合$\Gamma$是非空、紧值和连续的。
    因此，对任意$f\in C\left(X\right), x\in X$，如下问题：
    $$Tf\left(x\right)=\max_{y\in\Gamma\left(x\right)}\left[F\left(x,y\right)+\beta f\left(y\right)\right]$$
    等价于在紧集$\Gamma\left(x\right)$上最大化一个连续函数，从而最大值必能取到。同时，由于$F$和$f$均有界，从而$Tf$有界。
    最后，由最大化定理（定理\ref{th-max-th}），$Tf$（即定理中的函数$h$）连续。从而证明了：$T: C\left(X\right)\rightarrow C\left(X\right)$。

    \textbf{Step2：利用Blackwell条件，证明算子$T$是压缩映射}

    Blakwell充分条件（定理\ref{th-blackwell}）的前提：定义了上确界范数的有界函数空间，显然$C\left(X\right)$符合。
    接着验证两个条件
    \begin{itemize}
        \item 单调性：
        
            令$f,h \in C\left(X\right), f\leq h$，对$f$应用算子$T$：
            $$Tf\left(x\right)=\max_{y\in \Gamma\left(x\right)}\left[F\left(x,y\right)+\beta f\left(y\right)\right]$$
            对$\forall x\in X$，令$g\left(x;f\right)$为达到最大值的$y$的集合：
            $$
            g\left(x;f\right)\equiv \{y\in\Gamma\left(x\right)|Tf\left(x\right)=F\left(x,y\right)+\beta f\left(y\right)\}
            $$
            由于：$f\left(y\right)\leq h\left(y\right), \forall y\in \Gamma\left(x\right)$
            从而：$$Tf\left(x\right)=F\left(x,g\left(x;f\right)\right)+\beta f\left(g\left(x;f\right)\right)
            \leq F\left(x,g\left(x;f\right)\right)+\beta h\left(g\left(x;F\right)\right), \quad \forall x\in X$$
            注意到：
            $$
            F(x, g(x ; f))+\beta h(g(x ; f)) \leq \max _{y \in \Gamma(x)}[F(x, y)+\beta h(y)]=T h(x), \quad \forall x \in X
            $$
            因此有：$Tf\left(x\right)\leq Th\left(x\right), \quad x\in X$，即$T$满足单调性。
        \item discounting：
        
            对任意$f\in C\left(X\right)$：
            $$
            \begin{aligned}
                T\left(f+a\right)\left(x\right) &= 
                    \max_{y\in\Gamma\left(x\right)}\left[F\left(x,y\right)+\beta \left(f\left(y\right)+a\right)\right]\\
                    &= \max_{y\in\Gamma\left(x\right)}\left[F\left(x,y\right)+\beta \left(f\left(y\right)\right)\right]+\beta a
            \end{aligned}
            $$
            即：$T\left(f+a\right)\left(x\right)=T\left(x\right)+\beta a, \quad f\in C\left(X\right)$。
    \end{itemize}

    因此，算子$T$符合Blackwell充分条件，从而是$C\left(X\right)$上的压缩映射。

    \textbf{Step3：应用压缩映射定理}

    我们已经证明，定义了上确界范数的有界连续函数空间$C(X)$是完备空间。因此由压缩映射定理（定理~\ref{th-cmt}）立得后两条结论。
\end{proof}

定理\ref{th-fe-solution-exist}证明了：
\begin{itemize}
    \item 对于形一般化的泛函问题（$FE$）\ref{eq-general-fe}：
        $v\left(x\right)=\max_{y\in\Gamma\left(x\right)}\left[F\left(x,y\right)+\beta v\left(y\right)\right]$，不动点存在且唯一
    \item 数值解法的合理性：给定一个初始猜测，通过迭代算子$T$可以收敛到不动点
\end{itemize}

回顾定理\ref{th-solution-fe-sp}：满足特定有界条件的$FE$问题的解$v$就是$SP$问题的解$v^*$（上确界函数）。
根据这一最优性原理：在定理\ref{th-fe-solution-exist}的假设下，满足一般化$FE$问题式\ref{eq-general-fe}：
$v\left(x_t\right)=F\left(x_t, x_{t+1}\right)+\beta v\left(x_{t+1}\right), ~ t=0,1,2, \ldots$
的唯一有界连续函数$v$就是对应序列问题（$SP$）的上确界函数。也就是说，在假设\ref{as-convex-X}和假设\ref{as-bounded-returns}下，
定理\ref{th-solution-fe-sp}和定理\ref{th-fe-solution-exist}共同证明：\textbf{上确界函数有界并且连续}。

回顾定理\ref{th-policy-fe-sp}：满足
$v^*\left(x_t^*\right) = F\left(x_t^*,x_{t+1}^*\right) + \beta v^*\left(x_{t+1}^*\right), ~ t=0,1,2,\cdots$
和特定有界性条件的任意序列是一个最优计划。因此，结合定理\ref{th-policy-fe-sp}和定理\ref{th-fe-solution-exist}可知，至少存在一个最优计划：
由非空对应$G$产生的任意计划是最优的。


% 1.3.4.4 值函数性质
\subsubsection{值函数性质}



% 1.4 ===== 新古典增长模型初步数值解法 =====
\section{确定性最优增长模型的初步数值解法}

%
\subsection{值函数迭代}

%
\subsection{Euler Based Methods}

%
\subsection{Howard Improvment Method}

